\subsection{Test Case Design}
\label{sec:testcase_design}

The main text reports representative coverage only; the complete test-case specification, including inputs, expected results, and exception conditions, is provided in Appendix~\ref{appendix:testcases}.

\begin{table}[H]
\centering
\caption{Representative test coverage}
\label{tab:tc_representative}
\begin{tabularx}{\textwidth}{>{\raggedright\arraybackslash}p{0.23\textwidth} X X}
\hline
\textbf{Area} & \textbf{Representative coverage} & \textbf{Evidence of the expected behaviour} \\
\hline
Authentication and access control & Google SSO, MFA verification, and permission-protected routes & Successful authentication grants access; invalid second factors and insufficient permissions are rejected. \\
\hline
Learning loop and adaptive quiz & Quiz configuration, response recording, and SM-2 review updates & Published settings reach student attempts, response evidence is persisted, and the derived card state is updated. \\
\hline
Application loop and mock interview & Outcome-based interview configuration, adaptive turn-taking, and session completion & Eligible students can complete an interview and receive a persisted transcript, metadata, and verdict. \\
\hline
\end{tabularx}
\end{table}

These representatives demonstrate the three principal test levels used by the system: system behaviour at user-facing boundaries, integration behaviour across feature boundaries, and persistence of the resulting learning evidence. The appendix maps the individual cases to their functional requirements.
